\section{Related Work}

\paragraph{Point-cloud representation learning.}
The current codebase already touches core point-cloud backbones such as PointNet++ \citep{qi2017pointnetpp} and Point Transformer \citep{zhao2021pointtransformer}. A natural comparison set also includes Point-BERT \citep{yu2022pointbert} and PointNeXt \citep{qian2022pointnext}, which improve point-cloud representation quality through pretraining and stronger training recipes.

\paragraph{Single-view reconstruction and completion.}
The literature folder is already anchored in structure-aware reconstruction and completion, including StructureNet-style hierarchical generation and topology-aware reconstruction. Additional strong references include PoinTr \citep{yu2021pointr}, SnowflakeNet \citep{xiang2021snowflake}, and SeedFormer \citep{zhou2022seedformer}, all of which are relevant if the paper emphasizes partial-to-complete geometry reasoning.

\paragraph{Scene graphs and hierarchical spatial structure.}
The project is especially close to work that treats space as a structured graph rather than an unordered set. Relevant references include 3D semantic scene graph learning \citep{wald2022learning}, egocentric 3D scene graph prediction \citep{zhang2024egosg}, and multiview scene graphs \citep{zhang2024multiview}. These papers help position the localization stage as reasoning over structured scene memory rather than only pose regression.

\paragraph{Layout and hierarchical structure learning.}
The existing literature folder also includes hierarchical layout representation work \citep{bai2023layout} and indoor scene hierarchy generation. Those references are useful for motivating why multiscale structure matters when moving from local observations to global spatial organization.
